\documentclass[11pt]{article}

% --------------------
% Page setup
% --------------------
\usepackage[margin=1in]{geometry}
\usepackage{setspace}
\setstretch{1.1}

% --------------------
% Math and symbols
% --------------------
\usepackage{amsmath}
\usepackage{amssymb}

% --------------------
% Figures and tables
% --------------------
\usepackage{graphicx}
\usepackage{float}
\usepackage{booktabs}
\usepackage{caption}
\usepackage{subcaption}

% --------------------
% Lists and links
% --------------------
\usepackage{enumitem}
\usepackage[hidelinks]{hyperref}

% --------------------
% Code blocks
% --------------------
\usepackage{xcolor}
\usepackage{listings}

\lstset{
    basicstyle=\ttfamily\small,
    breaklines=true,
    frame=single,
    numbers=left,
    numberstyle=\tiny,
    tabsize=4
}

% --------------------
% Title
% --------------------
\title{\textbf{UIUC Basketball Project - Lineup Construction and Recruiting Dashboard}}
\author{Anirudhha S. Harsha}
\date{\today}

\begin{document}

\maketitle

\section{Problem Statment}

Transfer portal numbers have shot through the roof in recent years. For that reason, it has been exponentially more difficult to scout more than two thousand players every year as opposed to the ~900 guys that entered when the portal first started in 2018.

On top of that, GMs have to deal with scouting high schoolers, D2/D3 recruits, and international players. This amounts to a single college basketball scouting staff needing to sift through a pool of more than three or four thousand players for the next program-changing individual.

To help alleviate some of that stress, my website offers a transfer-portal and high-school centric scouting tool, which compiles transfer data and high school data from multiple sources. Various filters allow the user to search the pool for common metrics and needs as opposed to digging through hours of film. Furthermore, my website projects the value of high school players, transfers, and returning players through a statistic called Bayesian Performance Rating (developed by Evan Miyakawa). Finally, the "Simulator" tool on my website allows a GM to project their roster for next year while also receiving recommendations for high schoolers or transfers to add.

\section{Data}

The data I used were from the following sites:

\href{https://evanmiya.com}{Evan Miyakawa Basketball Page} - Used to get advanced stats on player for each season along with the important all-in-one metric BPR. Data was grabbed by writing a scraper to parse the website. This was more difficult, the website seems to be an R-shiny app structure.

\href{https://barttorvik.com/#}{Barttorvik} - Grabbed player stats for every year since 2010 (dunks, mid ranges were tracked after 2010). Also grabbed transfer data for every year since 2018. Barttorvik has a specific way to download the data as a .csv by searching up a URL in a specific structure, I did this to gather regular season player stats. I wrote a scraper to grab transfer stats since 2018, this info is not easily accessible using his website or the structured URLs.

\href{https://247sports.com/college/basketball/recruiting/}{247} - A classic for basketball scrapers. Used this to grab recruit information for all years since class of 2009 and all transfer portal information since 2018.

\href{https://collegebasketballportal.com/}{College Basketball Portal - Luke McCartney} - Not a source of data but offered me inspiration for my frontend and UX.


\section{Methodology}

\subsection{Name Matching}

The absolute worst part of working with sports data from multiple sources. The primary packages I used were \textbf{RapidFuzz} along with some manual matching with lower-confidence matches. This was a necessary step to stitch together data from the three sources I used, and unfortunately it did cause some loss in the initial data gathered from the three sites when training my ML models.

\subsubsection{Storing Data}

All of my data was stored using \textbf{DuckDB} as opposed to .csv files, though .csv files were used occasionally to help visualize data. Through the use of .db files, I had to make extensive use of \textbf{SQL} queries to both dig through my .db files and pull data into my models.

\subsubsection{Predicting High School/Transfer Playtype}

One thing I noticed on Barttorvik's site was that he had "Role" along with traditional position. I decided I wanted to predict "Role" as "Playtype" for HS recruits and transfers. I had categorical and numerical data for HS recruits as well as transfers in their previous year of college, so I used \textbf{CatBoost}, a tree-based model. In order to automatically tune CatBoost's hyperparameters I used \textbf{Optuna}.

The HS model was trained on 247 data and output player role. For HS recruits, the train/val/test split depended on the recruit year (train was 2024 and 2025), val was (2022 and 2023), and train was the rest of the years.

The transfer model was trained on previous-year Barttorvik numerical statistics and previous role and output their future role. Due to the skewed number of transfers based on the year, train/val/test was a 70/15/15 split between all rows available ordered chronologically.

The overall output of this model (whether for HS or for transfers) were more general playtypes such as Scoring PG, Pure PG, Wing F, Wing G, etc. that describe more of a player's tendencies than their 1-5 on-paper position.


\subsubsection{Predicting BPR}

One of the most important predictive metrics for the site was BPR. BPR was taken from Evan Miya's website and is a representation of a player's overall contribution to a team. Seeing as it is an all-encompassing stat such as WAR or WARP, I decided to predict it for HS players and transfers.

Since the data was more numerical in nature, I tested a variety of models for their performance: CatBoost, \textbf{XGBoost}, and \textbf{Ridge Classifiers}

For high schoolers, the input data was similar to the inputs for the playtype model. However, I also included more 247 ranking data since BPR is a performative metric while playtype is simply a classification. I also included the destination institution as a feature or input to the model, as some institutions may have a larger tendency to pump out higher-BPR athletes.

For high schoolers without a committed or enrolled institution, I trained the model without the destination. When predicting BPR with the 2026 class, I use either the model trained with destination or without destination depending on the status of the recruit.



\section{Results}

\section{Discussion}

\section{Conclusion}

\section*{References}

\appendix

\section{Appendix}

\end{document}